\section{The compatible cone}\label{sec:cones}

The distinction between existence and realization of a specified class
appears in the separating functional. For existence, it annihilates all
closed invariant forms. For a specified class, it need only annihilate
exact invariant forms. We first establish the estimates in this setting,
then identify the classes of the resulting currents.

\subsection{Invariant representatives and separation}

We retain the Hermitian metric, the projection \(R\), and the closed
lift \(B=dd^*G+H\) of Lemma~\ref{lem:lift}.
For a closed current \(D\), let
\(\mathfrak c(D)\in H^2(X;\R)\) denote its intersection-dual class:
\[
 D(\alpha)=\mathfrak c(D)\cdot[\alpha]
 \quad\text{for every smooth closed two-form }\alpha.
\]
Equivalently, \(\mathfrak c(D)\) is represented by the Hodge star of
the harmonic part of the distributional two-form associated with \(D\).

\begin{lemma}\label{lem:invariant-cohomology}
The classes represented by closed invariant forms are exactly
\(\mathcal V=(H_J^-)^\perp\). The intersection form is positive
definite on \(H_J^-\), and
\[
 H^2(X;\R)=\mathcal V\oplus H_J^-.
\]
The operator \(1-BR\) takes exact forms to exact invariant forms.
\end{lemma}

\begin{proof}
Closed anti-invariant forms are self-dual, hence harmonic. Their
intersection pairing is their \(L^2\) inner product, so it is positive
definite on their classes. Invariant and anti-invariant forms wedge to
zero pointwise.

For any closed two-form \(\alpha\), the form
\((1-BR)\alpha\) is closed and invariant, and
\[
 [\alpha]=[(1-BR)\alpha]+[HR\alpha],
\]
because \(BR\alpha=dd^*GR\alpha+HR\alpha\). The second summand
is represented by a closed anti-invariant form. The two subspaces are
orthogonal, and positive definiteness on the second makes their
intersection zero. This proves the decomposition and the assertion
about \(\mathcal V\).

If \(\alpha\) is exact, it is \(L^2\)-orthogonal to every harmonic
form. For a harmonic anti-invariant form \(\beta\),
\(\ip{R\alpha}{\beta}_{L^2}=\ip{\alpha}{\beta}_{L^2}=0\).
Thus \(HR\alpha=0\), and both \(BR\alpha\) and
\((1-BR)\alpha\) are exact.
\end{proof}

A current is called \emph{positive} here if it annihilates
anti-invariant forms and is nonnegative on pointwise semipositive
invariant forms. As in Proposition~\ref{prop:separation}, it has a
representation by a finite positive measure on \(\mathscr C\).
Its trace measure is the projection to \(X\).

\begin{lemma}[Prescribed-class separation]\label{lem:class-separation}
Let \(a\in\mathcal V\) have no compatible representative. There
are currents \(P,Q,T\in H^{-3}\), with \(P\) nonzero and positive,
such that
\begin{gather}
 P(F)=1,\qquad P\text{ annihilates exact invariant forms},
 \label{eq:class-exact-tests}\\
 T=P+Q\text{ is closed},\qquad RQ=Q=*Q,
 \label{eq:class-current}\\
 \mathfrak c(T)\in\mathcal V,\qquad
 \mathfrak c(T)\cdot a\le0.
 \label{eq:class-separator-sign}
\end{gather}
Moreover, a class \(b\in H^2(X;\R)\) is taming if and only if
\[
 b\cdot\mathfrak c(N)>0
 \quad\text{for every nonzero closed positive current }N.
\]
\end{lemma}

\begin{proof}
By Lemma~\ref{lem:invariant-cohomology}, the closed invariant
representatives of \(a\) form a nonempty affine space. Separate
this affine space from the open cone of positive invariant forms
in the Fr\'echet space of smooth invariant forms. Hahn--Banach
separation gives a nonzero continuous functional, nonnegative on
the positive cone and nonpositive on the affine space. It vanishes
on the direction space, which consists of exact invariant forms.
Indeed, translating by any real multiple of a direction preserves
the affine space, while positive scalings preserve the cone.
Extend the functional by \(1-R\) and normalize it as in
Proposition~\ref{prop:separation} to obtain \(P(F)=1\).
The same positivity argument gives its measure representation.

Define \(T(\alpha)=P((1-BR)\alpha)\) and
\(Q(\alpha)=-P(BR\alpha)\). Lemma~\ref{lem:invariant-cohomology}
shows that \(T\) annihilates every exact form, so it is closed.
The projection identities give \(RQ=Q=*Q\). A closed
anti-invariant form \(\beta\) lies in the kernel of the operator
used to construct \(G\); hence \(B\beta=\beta\) and
\(T(\beta)=0\). Thus \(\mathfrak c(T)\in\mathcal V\).
On a closed invariant representative of \(a\), the currents
\(T\) and \(P\) agree, proving \eqref{eq:class-separator-sign}.
The Sobolev assertion follows from the measure representation and
the order-zero lift, as in Proposition~\ref{prop:separation}.

For the last statement, a taming representative evaluates strictly
positively on every nonzero positive current. Conversely, if \(b\)
has no taming representative, separate its affine space of closed
representatives from the open cone of forms positive on complex
lines. The separator annihilates all exact forms. Since arbitrary
anti-invariant forms can be added to that cone, it also annihilates
all anti-invariant forms. It is therefore a nonzero closed positive
current \(N\), and \(b\cdot\mathfrak c(N)\le0\).
\end{proof}

\subsection{Pairings for currents of arbitrary class}

\begin{lemma}\label{lem:cohomological-pairing}
For closed currents \(D_1,D_2\in H^{-3}\),
\begin{equation}\label{eq:cohomological-pairing}
 \ip{S_rD_1}{*S_rD_2}_{L^2}
 =\mathfrak c(D_1)\cdot\mathfrak c(D_2).
\end{equation}
\end{lemma}

\begin{proof}
The forms \(*S_rD_i\) lie in \(H^3\), are distributionally closed,
and have the harmonic components representing \(\mathfrak c(D_i)\).
Here we use commutation with \(d^*,*\), and the fact that \(S_r\)
is the identity on harmonic forms. Their heat regularizations are
smooth closed representatives of these classes. The intersection
identity for smooth forms passes to the limit in \(L^2\).
In degree two on a four-manifold,
\(\alpha\wedge\beta=(*\alpha)\wedge(*\beta)\), giving
\eqref{eq:cohomological-pairing}.
\end{proof}

We record the estimates that apply to the separator in
Lemma~\ref{lem:class-separation}. In the next proposition, positive
currents are not required to have any prescribed cohomology class.

\begin{proposition}\label{prop:class-estimates}
Suppose that \(P\ne0\) is positive, annihilates exact invariant
forms, and \(T=P+Q\) is closed, where \(Q\in H^{-3}\) is
anti-invariant. Write \(\lambda\) for a positive measure representing
\(P\), and \(\mu\) for its trace. Then
\begin{gather}
 \mu(B_s(x))\le Cs^2,\qquad Q\in L^2,
 \label{eq:class-growth}\\
 \iint_{d_g(x,y)<r}|\ell-\tau_{yx}m|^2
       \dd\lambda(x,\ell)\dd\lambda(y,m)\le Cr^4.
 \label{eq:class-angular}
\end{gather}
Consequently the positive-density restriction \(P_E\) is closed,
where \(E=\{x:\lim_{s\downarrow0}s^{-2}\mu(B_s(x))>0\}\).

If \(A\) is any positive current whose trace has quadratic growth,
then
\begin{equation}\label{eq:class-mixed-limit}
 \lim_{r\downarrow0}\ip{S_rA}{*S_rQ}_{L^2}=0.
\end{equation}
Every closed positive current has quadratic trace growth.
\end{proposition}

\begin{proof}
Each radial test \(Du=(1-BR)dd_J^cu\) in
Section~\ref{sec:mass} is exact by
Lemma~\ref{lem:invariant-cohomology}. Its pointwise estimates
therefore give the quadratic growth bound under the present
annihilation hypothesis. The normalization \(P(F)=1\) is inessential:
for arbitrary finite mass it only changes the constant. In particular,
this argument applies to every closed positive current.

We will use the kernel estimates for any pair \(A_1,A_2\) of positive
currents with quadratic-growth trace measures \(\mu_1,\mu_2\).
Write \(\lambda_1,\lambda_2\) for their representing measures.
The proof of Lemma~\ref{lem:positive-pairing}, using
Lemma~\ref{lem:kernel} and \eqref{eq:angular-identity}, gives
\begin{align}
 \|S_rA_i\|_2&\le C_i r^{-1},\label{eq:general-smoothed-mass}\\
 \ip{S_rA_1}{*S_rA_2}_{L^2}
 &\ge cr^{-4}\iint_{d<r}|\ell-\tau_{yx}m|^2
            \dd\lambda_1(x,\ell)\dd\lambda_2(y,m)
 \nonumber\\
 &\quad-Cr^{-2}\iint(1+d/r)^{-6}
                     \dd\mu_1(x)\dd\mu_2(y),
 \label{eq:general-positive-pairing}
\end{align}
where \(d=d_g(x,y)\). These statements do not require either measure
to be dominated by the other. The shell bound
\eqref{eq:shell} holds separately for each \(\mu_i\), proving
the norm bound and a uniform bound on the negative term.

For any anti-invariant distribution \(W\) with \(S_rW\in L^2\),
Lemma~\ref{lem:mixing} and \eqref{eq:general-smoothed-mass} give
\begin{equation}\label{eq:general-mixed-bound}
 |\ip{S_rA_i}{*S_rW}_{L^2}|
 \le Cr\|S_rA_i\|_2\|S_rW\|_2
 \le C_i'\|S_rW\|_2.
\end{equation}
Apply \eqref{eq:cohomological-pairing} to \(T,T\).
Since \(Q\) is self-dual, expansion and these estimates yield
\[
 cr^{-4}\ang_r+\|S_rQ\|_2^2
 \le C+|\mathfrak c(T)^2|+C\|S_rQ\|_2,
\]
where \(\ang_r\) is the angular integral in
\eqref{eq:class-angular}. Hence \(S_rQ\) is uniformly bounded
in \(L^2\), and the angular bound follows. Weak compactness and
distributional convergence identify an \(L^2\) limit with \(Q\).

For smooth anti-invariant \(W\),
\[
 \ip{S_rA}{*S_rW}=A(S_r^2W)\longrightarrow A(W)=0.
\]
The uniform bound \eqref{eq:general-mixed-bound}, the \(L^2\)
contraction property of \(S_r\), and approximation by smooth
anti-invariant forms give \eqref{eq:class-mixed-limit}.
Finally, \eqref{eq:class-growth}--\eqref{eq:class-angular} verify
the hypotheses of Theorem~\ref{thm:splitting}, which proves that
\(P_E\) is closed.
\end{proof}

The next lemma permits pairing the zero-density part with positive
currents unrelated to the original separator.

\begin{lemma}[Residual pairings]\label{lem:general-residual}
Let \(A,P'\) be positive currents whose trace measures \(\nu,\mu'\)
have quadratic growth. Suppose
\[
 \lim_{s\downarrow0}s^{-2}\mu'(B_s(x))=0
 \quad\text{for }\mu'\text{-almost every }x.
\]
Then
\begin{equation}\label{eq:general-residual}
 \liminf_{r\downarrow0}\ip{S_rA}{*S_rP'}_{L^2}\ge0.
\end{equation}
\end{lemma}

\begin{proof}
We first show
\begin{equation}\label{eq:cross-small-density}
 (\mu'\times\nu)\{(x,y):d_g(x,y)<s\}=o(s^2).
\end{equation}
For small \(s\), choose an \(s\)-separated maximal set of centers
and let \(B_j\) be their radius-\(s\) balls. They cover \(X\),
and their fixed dilates have uniformly bounded overlap, by the
volume comparison in smooth coordinate charts. If \(x\in B_j\)
and \(d_g(x,y)<s\), then \(y\in2B_j\). Thus
\[
 (\mu'\times\nu)\{d<s\}
 \le\sum_j\mu'(B_j)\nu(2B_j).
\]
Bounded overlap gives
\[
 \sum_j\mu'(B_j)^2
 \le C\int_X\mu'(B_{2s}(x))\dd\mu'(x)=o(s^2).
\]
Indeed, the integrand divided by \(s^2\) is bounded by quadratic
growth and tends to zero at \(\mu'\)-almost every center.
Also
\[
 \sum_j\nu(2B_j)^2
 \le Cs^2\sum_j\nu(2B_j)=O(s^2).
\]
Cauchy--Schwarz proves \eqref{eq:cross-small-density}.

In \eqref{eq:general-positive-pairing}, discard the angular term.
The remaining negative error tends to zero: on \(d<r\), and on
each fixed shell \(2^jr\le d<2^{j+1}r\), this follows from
\eqref{eq:cross-small-density}. Quadratic growth bounds the shell
contributions by \(C2^{-4j}\), allowing summation by dominated
convergence. This proves \eqref{eq:general-residual}.
\end{proof}

\subsection{Proof of the cone description}

\begin{proof}[Proof of Theorem~\ref{thm:taming-cones}]
Adding a closed anti-invariant form does not change positivity on
complex lines. Lemma~\ref{lem:invariant-cohomology} therefore gives
\eqref{eq:tame-decomposition}. In particular \(\mathcal C\) is
nonempty. It is a convex cone of positive-square classes in
\(\mathcal V\), since its classes have closed taming representatives.

A compatible form wedges strictly positively with every taming form.
Indeed, in a unitary frame diagonalizing the compatible form, the
wedge product is a positive linear combination of the two positive
diagonal entries of the taming form's invariant part. Consequently
every \(a\in\mathcal K_J^c\) lies in \(\mathcal C\) and pairs
positively with every class in \(\mathcal C\), hence nonnegatively
with \(\overline{\mathcal C}\).

The compatible cone is open in \(\mathcal V\): add sufficiently
small linear combinations of closed invariant representatives of a
basis of \(\mathcal V\) to a compatible form. The intersection form
on \(\mathcal V\) is nondegenerate by
Lemma~\ref{lem:invariant-cohomology}. If \(a\cdot y=0\) for a
nonzero \(y\in\overline{\mathcal C}\), a sufficiently small
perturbation of \(a\) within the compatible cone would pair
negatively with \(y\). Thus every compatible class satisfies the
strict positivity in \eqref{eq:compatible-duality}.

Conversely, suppose \(a\in\mathcal C\) has this strict positivity
but no compatible representative. Apply
Lemma~\ref{lem:class-separation} and
Proposition~\ref{prop:class-estimates}. Let
\[
 P'=P-P_E,\qquad T'=P'+Q=T-P_E,\qquad
 y=\mathfrak c(T').
\]
Both \(P_E\) and \(T'\) are closed. Since \(P_E\) is invariant,
\(\mathfrak c(P_E)\in\mathcal V\), and hence \(y\in\mathcal V\).
The trace \(\mu'\) of \(P'\) satisfies the hypotheses of
Lemma~\ref{lem:general-residual}: it is bounded by \(\mu\) and
is concentrated where the latter has zero two-density.

We claim that \(y\in\overline{\mathcal C}\). Let \(N\) be any
nonzero closed positive current. Its trace has quadratic growth by
Proposition~\ref{prop:class-estimates}. By
\eqref{eq:cohomological-pairing},
\[
 y\cdot\mathfrak c(N)
 =\ip{S_rN}{*S_rP'}+\ip{S_rN}{*S_rQ}.
\]
The mixed term tends to zero by \eqref{eq:class-mixed-limit}, and
Lemma~\ref{lem:general-residual} bounds the lower limit of the first
term by zero. Thus \(y\cdot\mathfrak c(N)\ge0\). Since \(a\)
is taming, \((y+\eps a)\cdot\mathfrak c(N)>0\) for every
\(\eps>0\). The last part of Lemma~\ref{lem:class-separation}
shows that \(y+\eps a\) is taming. It also lies in \(\mathcal V\),
so \(y+\eps a\in\mathcal C\), proving the claim.

It follows that \(a\cdot y\ge0\), with strict inequality if
\(y\ne0\). If \(P'\ne0\), then \(y\ne0\). To see this,
suppose \(y=0\) and pair \(T'\) with itself using
\eqref{eq:cohomological-pairing}. We obtain
\[
 0=\ip{S_rP'}{*S_rP'}+2\ip{S_rP'}{*S_rQ}+\|S_rQ\|_2^2.
\]
The residual lemma, the mixed-term limit, and strong \(L^2\)
convergence of \(S_rQ\) force \(Q=0\). Then the nonzero positive
current \(P'=T'\) evaluates strictly positively on a closed tamer,
contrary to \(\mathfrak c(T')=0\).

Finally, a taming representative of \(a\) gives
\(a\cdot\mathfrak c(P_E)\ge0\), strictly if \(P_E\ne0\).
At least one of \(P_E,P'\) is nonzero. Therefore
\[
 a\cdot\mathfrak c(T)
 =a\cdot\mathfrak c(P_E)+a\cdot y>0,
\]
contradicting \eqref{eq:class-separator-sign}. This proves
\eqref{eq:compatible-duality}.
\end{proof}

\begin{proof}[Proof of Corollary~\ref{cor:taming-cones}]
The subspace \(H_J^-\) is positive definite, while its orthogonal
complement \(\mathcal V\) contains the nonempty positive-square
cone \(\mathcal C\). Hence \(h_J^-\le b_2^+-1\). If equality
holds, the intersection form on \(\mathcal V\) has signature
\((1,b_2^-)\). Convexity puts \(\mathcal C\) in one component
of its positive-square cone. Every element of that component pairs
strictly positively with every nonzero element of its closure,
by Cauchy--Schwarz in coordinates of signature
\((+,-,\ldots,-)\). Thus all \(a\in\mathcal C\) satisfy
\eqref{eq:compatible-duality}, and
\(\mathcal C=\mathcal K_J^c\). Equation
\eqref{eq:tame-decomposition} gives the asserted identity.
If \(b_2^+=1\), the inequality forces \(h_J^-=0\), proving the
last assertion.
\end{proof}

\begin{example}[Strict inclusion on the four-torus]\label{ex:strict-cones}
On \(X=\R^4/\mathbb Z^4\), use the Euclidean metric and standard
orientation, and put \(e^{ij}=dx_i\wedge dx_j\). Set
\[
 U=e^{12}+e^{34},\qquad V_0=e^{13}-e^{24},\qquad
 W=e^{14}+e^{23},
\]
\[
 f=\tfrac14\sin(2\pi x_1),\qquad
 k=\tfrac14\sin(2\pi x_2),\qquad
 r=(1+f^2+k^2)^{1/2}.
\]
The form \(F=(U+fV_0+kW)/r\) is self-dual and has squared norm
two, so it is the fundamental form of a smooth orthogonal almost
complex structure \(J\) of this orientation. The form
\[
 \eta=U+2f e^{13}+2k e^{23}
\]
is closed. Its self-dual part is \(rF\), so it is invariant, and
its anti-self-dual part is
\[
 f(e^{13}+e^{24})+k(e^{23}-e^{14}).
\]
Its Hermitian eigenvalues are \(r\pm\sqrt{f^2+k^2}>0\).
Thus \(\eta\) is compatible. The sine terms are exact, giving
\([\eta]=[U]\).

A closed anti-invariant form is self-dual and harmonic for the flat
metric, hence is a constant combination \(aU+bV_0+cW\).
Orthogonality to \(F\) gives \(a+bf+ck=0\) everywhere. Varying
\(x_1,x_2\) yields \(a=b=c=0\); thus \(H_J^-=0\).
The closed forms \(U\pm V_0\) have invariant projections
\((1\pm f)F/r\), so both tame \(J\). Their classes pair to zero,
since \(U^2=V_0^2\). Neither class can have a compatible
representative, which would wedge strictly positively with a tamer
in the other class. Hence
\[
 \mathcal K_J^c+H_J^-=\mathcal K_J^c
 \subsetneq\mathcal K_J^t.
\]
In Theorem~\ref{thm:taming-cones}, the class \([U+V_0]\) fails
the strict positivity test against \([U-V_0]\in\mathcal C\),
and conversely.
\end{example}
